\begin{figure}[t]
\centering
\begin{tikzpicture}[
  font=\scriptsize,
  box/.style={draw, rounded corners=2pt, align=center, inner sep=3pt,
              text width=54mm, minimum height=6mm, line width=0.5pt},
  keep/.style={box, draw=cIdentity, fill=cIdentity!8},
  narrow/.style={box, draw=cNeutral, fill=cFaint},
  final/.style={box, draw=cAccent, fill=cAccent!8},
  down/.style={-{Stealth[length=1.8mm]}, line width=0.6pt},
  side/.style={font=\tiny, align=left, text width=20mm, inner sep=1pt},
]

\node[keep] (a) {\textbf{\nReposCollected} public repositories collected, no
  query hit the API ceiling};
\node[narrow, below=4mm of a] (b) {\textbf{\nReposWithRuntime} of them declare
  an agent runtime, holding \textbf{\nTemplateDirs} template directories};
\node[narrow, below=4mm of b] (c) {after removing \nTemplateExcluded test
  fixtures, snapshots and vendored copies, and \nTemplateForkCopies fork
  copies};
\node[final, below=4mm of c] (d) {\textbf{\nProjects} analysed Terraform
  projects\\\scriptsize denominator for escalation results};
\node[final, below=13mm of d] (e) {\textbf{\nChainDenominator} units covered by
  the chain analysis\\\scriptsize \nChainTerraform Terraform $+$ \nChainCfn
  CloudFormation};

\draw[down] (a) -- (b);
\draw[down] (b) -- (c);
\draw[down] (c) -- (d);
\draw[down] (d) -- node[side, right=1mm]
  {\nChainNotCovered units not covered, named individually} (e);

\end{tikzpicture}
\caption{How the collection narrows to each denominator. The two shaded boxes
are different populations and the paper does not mix them: a rate over
\nProjects and a rate over \nChainDenominator answer different questions.}
\label{fig:pipeline}
\end{figure}
